\SourceLeaf{347}%
\par\indent Bei Untersuchungen, wo eine Mannigfaltigkeit von Richtungen gerader Linien im Raume ins Spiel kommt, ist es vortheilhaft, diese Richtungen durch diejenigen Punkte auf der Oberfläche einer festen Kugel zu bezeichnen, welche die Endpunkte der mit jenen parallel gezogenen Radien sind: Mittelpunkt und Halbmesser dieser \textit{Hülfskugel} sind hierbei ganz willkürlich; für letztern mag die Lineareinheit gewählt werden. Diess Verfahren kommt im Grunde mit demjenigen überein, welches in der Astronomie in stetem Gebrauch ist, wo man alle Richtungen auf eine fingirte Himmelskugel von unendlich grosem Halbmesser bezieht. Die sphärische Trigonometrie, und einige andere Lehrsätze, welchen der Verf. noch einen neuen von häufiger Anwendbarkeit beigefügt hat, dienen dann zur Auflösung der Aufgaben, welche die Vergleichung der verschiedenen vorkommenden Richtungen darbieten kann.

Wenn man die Richtung der an jedem Punkt einer krummen Fläche auf diese errichteten Normale durch den nach dem angedeuteten Verfahren entsprechenden Punkt der Kugelfläche bezeichnet, also jedem Punkt der krummen Fläche in dieser Beziehung einen Punkt der Oberfläche der Hülfskugel entsprechen lässt, so wird, allgemein zu reden, jeder Linie auf der krummen Fläche eine Linie auf der Oberfläche der Hülfskugel, und jedem Flächenstück von jener ein Flächenstück von dieser entsprechen. Je geringer die Abweichung jenes Stücks von der Ebene ist, desto kleiner wird der entsprechende Theil der Kugelfläche sein, und es ist mithin ein sehr natürlicher Gedanke zum Maassstabe der Totalkrümmung, welche einem Stück der krummen Fläche beizulegen ist, den Inhalt des entsprechenden Stücks der Kugelfläche zu gebrauchen. Der Verf. nennt daher diesen Inhalt die \textit{ganze Krümmung} des entsprechenden Stücks der krummen Fläche. Ausser der Grösse kommt aber zugleich noch die \textit{Lage} der Theile in Betracht, die, ganz abgesehen von dem Grössenverhältniss, in den beiden Stücken entweder eine ähnliche, oder eine verkehrte sein kann: diese beiden Fälle werden durch das der Totalkrümmung vorzusetzende positive oder negative Zeichen unterschieden werden können. Diese Unterscheidung hat jedoch nur insofern eine bestimmte Bedeutung, als die Figuren auf bestimmten Seiten der beiden Flächen gedacht werden: der Verf. nimmt sie bei der Kugelfläche auf der äussern und bei der krummen Fläche auf derjenigen Seite, wo man sich die Normale errichtet denkt, und es folgt dann, dass das positive Zeichen bei convex-convexen oder concav-concaven Flächen (die nicht wesentlich verschieden sind), und das nega%
